\documentclass[12pt]{book}               % I'm using a double-sided book style
\usepackage{amsmath}
\usepackage{amsfonts}
\usepackage{amsbsy}
\usepackage[sectionbib]{natbib}          % Cross-reference package (Natural BiB)
\usepackage{fancyhdr}                    % Fancy Header and Footer
\usepackage{setspace}                    % Line spacing
\usepackage{txfonts}                     % Public Times New Roman text & math font
\usepackage{graphicx}
\usepackage{hyperref}

\hypersetup{
    colorlinks=true,
    linkcolor=blue,
    filecolor=magenta,      
    urlcolor=cyan,
    citecolor=blue,
    }

%%% Fancy Header %%%%%%%%%%%%%%%%%%%%%%%%%%%%%%%%%%%%%%%%%%%%%%%%%%%%%%%%%%%%%%%%%%
% Fancy Header Style Options
\pagestyle{fancy}                       % Sets fancy header and footer
\fancyfoot{}                            % Delete current footer settings
\renewcommand{\chaptermark}[1]{         % Lower Case Chapter marker style
    \markboth{\chaptername\ \thechapter.\ #1}{}} %
\renewcommand{\sectionmark}[1]{         % Lower case Section marker style
    \markright{\thesection.\ #1}}         %
\fancyhead[LE,RO]{\bfseries\thepage}    % Page number (boldface) in left on even
                                        % pages and right on odd pages
\fancyhead[RE]{\bfseries\leftmark}      % Chapter in the right on even pages
\fancyhead[LO]{\bfseries\rightmark}     % Section in the left on odd pages
\renewcommand{\headrulewidth}{0.3pt}    % Width of head rule
%%% Clear Header %%%%%%%%%%%%%%%%%%%%%%%%%%%%%%%%%%%%%%%%%%%%%%%%%%%%%%%%%%%%%%%%%%
% Clear Header Style on the Last Empty Odd pages
\makeatletter
\def\cleardoublepage{\clearpage\if@twoside \ifodd\c@page\else%
    \hbox{}%
    \thispagestyle{empty}%              % Empty header styles
    \newpage%
    \if@twocolumn\hbox{}\newpage\fi\fi\fi}
\makeatother
\title{\includegraphics[width=\textwidth]{images/redbird_banner.png}
-- a comprehensive model-based image reconstruction platform for diffuse optical tomography and microwave tomography}


\author{Qianqian Fang}

\date{Version: 0.5}


\begin{document}

\maketitle              % you need to define \title{..}

\tableofcontents

\chapter{History}

Redbird is an open-source MATLAB/GNU Octave toolbox designed for solving the forward and inverse
problems for diffuse optical tomography (DOT), with the additional support of 
several other non-linear model-based imaging modalities. Redbird was ported from 
Redbird90~\cite{Fang2008} -- a FORTRAN90 based DOT image reconstruction software written by the same author, Dr. Qianqian Fang, between 2005 and 2008. The FORTRAN90-based Redbird90 was derived from an earlier implementation, code-named ``WhiteDragon''~\cite{Fang2004, Fang2010c}, also developed by Dr. Fang, between 2001 and 2005 as part of his Ph.D. dissertation~\cite{FangThesis}, designed to perform iterative 3-D microwave tomography (MWT) image reconstructions. The majority of the functionalities in WhiteDragon, including the 3-D finite element method (FEM) for forward modeling~\cite{Fang2004}, Gauss-Newton based nonlinear inverse solver~\cite{Fang2009IEEE}, a multi-right-hand-side quasi-minimal residual (BL-QMR)~\cite{Boyse1996} sparse linear solver, the adjoint method~\cite{Fang2004} for building Jacobians, and a dual-mesh reconstruction approach~\cite{Paulsen1995}, had been inherited to Redbird90, and subsequently ported to Redbird. To extend the modality from MWT to DOT, Redbird further introduced multi-spectral DOT reconstruction, derived from the multi-spectral MWT reconstruction algorithm developed in Dr. Fang's thesis~\cite{FangThesis}, a compositional-prior-guided DOT reconstruction~\cite{Fang2010,Deng2015a}, as well as a simultaneous source-detector (SD) coupling coefficient estimation~\cite{Boas2001,Schweiger2007}. 

Both names ``Redbird'' and ``WhiteDragon'' were rooted in ancient Chinese astrology. In ancient China, astrologists partitioned the sky into 4 quadrants -- each named after an animal deity that the constellations resemble: East -- pale/white dragon, West -- white tiger, North -- black turtle, and South -- red bird. The name ``Redbird'' was selected because DOT uses red and near-infrared light for imaging.

This document summarizes the mathematical treatment for the forward 
and inverse problems used in this algorithm.

\chapter{Mathematical models}

\section{DOT forward model}

The diffusion equation in the time-domain can be expressed as~\cite{Arridge1999}
\begin{equation}
\label{eq:diffusionTD} -\nabla\cdot
D(r)\nabla\Phi(r,t)+\mu_a(r)\Phi(r,t)+\frac{1}{c}\frac{\partial
\Phi(r,t)}{\partial t}=S_0(r,t)
\end{equation}
where $D(r)=\frac{1}{3(\mu_a+\mu_s')}$ is the diffusion
coefficient (unit cm); $c=\frac{c_0}{n}$ is the speed of light in the
medium (unit cm/s); $S_0(r,t)$ is the source. With assumed time dependence
$\exp(j\omega t)$, the frequency-domain (FD) diffusion equation can be
written as
\begin{equation}
\label{eq:diffusionFD} -\nabla\cdot
D(r)\nabla\Phi(r)+\left(\mu_a(r)+\frac{j\omega}{c}\right)\Phi(r)=S_0(r)
\end{equation}
where $\omega$ is the angular frequency and $\Phi(r)$ is the
phasor of $\Phi(r,t)$.

Integrating both sides of (\ref{eq:diffusionFD}) with a set of weight functions $\varphi_j(r) (j=1,...,M)$
where $M$ is the total number of basis functions, over the forward space $\Omega$, subsequently, applying the following vector
identity
\begin{equation}
\label{eq:veciden}\nabla\cdot f(r)\vec{g}(r)=f(r)\nabla\cdot
\vec{g}(r)+\vec{g}(r)\cdot\nabla f(r),
\end{equation}
we get
\begin{eqnarray}
\label{eq:femstep1}
-\int_{\Omega}{(D(r)\nabla^2\Phi(r))\varphi_j(r)dr}&-&\int_{\Omega}{\nabla D(r)\cdot\nabla\Phi(r)
\varphi_j(r)dr}\\\nonumber &+&\int_{\Omega}
{\left(\mu_a(r)+\frac{j\omega}{c}\right)\varphi_j(r)\Phi(r)dr}=\int_{\Omega}
{S_0(r)\varphi_j(r)dr}.
\end{eqnarray}
Assume $D(r)$ is constant for each forward element, i.e. element-based properties, 
we have $\nabla D(r)=0$ and the second term in (\ref{eq:femstep1}) becomes zero. 
However, if we define the optical properties on the nodes using a reconstruction mesh,
the second term will be non-zero. In this case, we expand $D(r)$ and $\mu_{a}(r)$ by
piecewise linear basis functions
\begin{equation}
\label{eq:parammesh}D(r)=\sum_{i}D_i\phi_i(r) \\
\mu_{a}(r)=\sum_{i}\mu_{a_i}\phi_i(r)
\end{equation}

as well as the Green's first identity
\begin{equation}
\label{eq:Green1} \int_{\Omega}{u(r)\nabla^2
v(r)dr}=-\int_{\Omega}{\nabla u(r)\cdot\nabla
v(r)dr}+\int_{\partial\Omega} {u(r)\nabla v(r)dr}
\end{equation}
we can rewrite (\ref{eq:diffusionFD}) as
\begin{eqnarray}
\label{eq:fem}
\int_{\Omega}{D(r)\nabla\varphi_j(r)\cdot\nabla\Phi(r)dr}&-&\int_{\partial\Omega}{D(r)\varphi_j(r)\nabla\Phi(r)\cdot
d\hat{r}}\\\nonumber &+&\int_{\Omega}
{\left(\mu_a(r)+\frac{j\omega}{c}\right)\varphi_j(r)\Phi(r)dr}=\int_{\Omega}
{S_0(r)\varphi_j(r)dr}
\end{eqnarray}
With Galerkin's method, i.e. the basis function is identical as
the weight function, $\Phi(r)$ is expanded as
$\Phi(r)=\sum_{i=1}^{4}\Phi_i\varphi_i(r)$ over each linear
Largarange forward element, and Equ. (\ref{eq:fem}) becomes
\begin{eqnarray}
\label{eq:galerkin}\nonumber&&\sum_{i}\Phi_i\left(\left\langle
D(r)\nabla\varphi_i(r)\cdot\nabla\varphi_j(r)\right\rangle -
\left\langle
D(r)\nabla\varphi_i(r)\varphi_j(r)\right\rangle_{\partial\Omega}\right.
\\&&+ \left.\left\langle
\left(\mu_a(r)+
\frac{j\omega}{c}\right)\varphi_i(r)\varphi_j(r)\right\rangle\right)=\left\langle
S_0(r)\varphi_j(r)\right\rangle
\end{eqnarray}
where $\langle u(r)\rangle$ denotes $\int_{\Omega}u(r)dr$ and
$\langle u(r)\rangle_{\partial\Omega}$ denotes
$\int_{\partial\Omega}u(r) dr$.

Expanding $\mu_a(r)$ and $D(r)$ on the parameter mesh basis as
$\mu_a(r)=\sum_{k=1}^{4}\mu_a^k\phi_k(r)$ and
$D(r)=\sum_{k=1}^{4}D_k\phi_k(r)$, we get the Galerkin weak form
equation as
\begin{eqnarray}
\label{eq:weakform}\nonumber&&\sum_{i}\Phi_i\left(\sum_{k}D_k\left\langle
\phi_k(r)\nabla\varphi_i(r)\cdot\nabla\varphi_j(r)\right\rangle -
\sum_{k}D_k\left\langle
\phi_k(r)\nabla\varphi_i(r)\varphi_j(r)\right\rangle_{\partial\Omega}\right.
\\&&\left.+\sum_{k}\left(\mu_{a}^k+\frac{j\omega}{c}\right)\left\langle\phi_k(r)\varphi_i(r)\varphi_j(r)\right\rangle\right)=\left\langle
S_0(r)\varphi_j(r)\right\rangle
\end{eqnarray}
where $\phi_k$ is the basis function of the parameter mesh.

\section{Boundary condition}

If the extrapolation boundary condition is used, the boundary
integration term in (\ref{eq:weakform}) is then dropped out,
leaving
\begin{eqnarray}
\label{eq:weakformzerobd}\sum_{i}&&\Phi_i\left(\sum_{k}D_k\left\langle
\phi_k(r)\nabla\varphi_i(r)\cdot\nabla\varphi_j(r)\right\rangle\right.\\\nonumber
+&&\left.\sum_{k}\left(\mu_{a}^k+\frac{j\omega}{c}\right)\left\langle\phi_k(r)\varphi_i(r)\varphi_j(r)\right\rangle\right)=\left\langle
S_0(r)\varphi_j(r)\right\rangle
\end{eqnarray}

However, it has been shown in Haskell \emph{et al.}~\cite{Haskell1994}, the partial current boundary
condition is more accurate. In this boundary condition, the fluence
satisfies the following condition on the boundary $\partial\Omega$

\begin{equation}
\label{eq:partialcurrent} \left(\boldsymbol\Phi=l_s\nabla{\boldsymbol\Phi}\right)|_{\partial\Omega}
\end{equation}
where $l_s$ is defined in Eq. 2.4.1 in Haskell, as
\begin{equation}
l_s=\frac{1+R_{eff}}{1-R_{eff}}2D
\end{equation}
and $R_{eff}=\frac{R_\phi+R_j}{2-R_\phi+R_j}$ is the effective reflection coefficient, and 
\begin{eqnarray}
R_\phi &=& \int_0^{\pi/2}2\sin \theta \cos\theta R_{Fresnel(\theta)}d\theta\\
R_j &=& \int_0^{\pi/2}3\sin \theta \cos^2\theta R_{Fresnel(\theta)}d\theta \\
R_{Fresnel}(\theta) &=& \frac{1}{2}\left(\frac{n \cos\theta' - n_{out} \cos\theta}{n \cos\theta' + n_{out} \cos\theta}\right)^2 \\
   &+&\frac{1}{2}\left(\frac{n \cos\theta - n_{out} \cos\theta'}{n \cos\theta + n_{out} \cos\theta'}\right)^2
\end{eqnarray}

Replacing $\nabla{\boldsymbol\Phi}=\boldsymbol\Phi/l_s$ to the boundary term in (\ref{eq:weakform}), we
have

\begin{eqnarray}
\label{eq:weakformpartial}\nonumber&&\sum_{i}\Phi_i\left(\sum_{k}D_k\left\langle
\phi_k(r)\nabla\varphi_i(r)\cdot\nabla\varphi_j(r)\right\rangle -
\sum_{k}\frac{1-R_{eff}}{2(1+R_{eff})}\left\langle
\phi_k(r)\varphi_j(r)\right\rangle_{\partial\Omega}\right.
\\&&\left.+\sum_{k}\left(\mu_{a}^k+\frac{j\omega}{c}\right)\left\langle\phi_k(r)\varphi_i(r)\varphi_j(r)\right\rangle\right)=\left\langle
S_0(r)\varphi_j(r)\right\rangle
\end{eqnarray}
and the boundary integration $\langle\phi_k(r)\varphi_j(r)\rangle_{\partial\Omega}=\frac{A}{6}$ if $k=j$ 
or $\frac{A}{12}$ if $k\ne j$. The general expression for the surface element integration is
\begin{equation}
\int_{\Omega_e}\varphi_1^l\varphi_2^m\varphi_3^ndr=\frac{l!m!n!}{(l+m+n+2)!}2A_e
\end{equation}


Equation (\ref{eq:weakformpartial}) is the final formula used in
the reconstruction code.

\section{Widefield illumination}
Unlike MC methods, DE is not capable of modeling propagation of collimated beams inside non-scattering medium such as air and water. This can be a challenge when simulating a light source that is located outside of the domain, i.e. a wide-field source.

In Redbird, a widefield source is converted to spatially distributed isotropic sources along the external surface at which the source is projecting. For example, when a disk source is casted on the top surface of a cubic diffusive domain projecting downward, a virtual source plane is created $1/\mu_s'$ below the actual domain boundary. Along this virtual source plane, all tetrahedral edges that intersect with this plane are identified. The source illumination aperture - i.e. the area where photons are launched, is subsequently sunken along this virtual source plane, and intersects all tetrahedral elements to break into smaller illumination areas bounded by each intersecting tetrahedron. The areas of these intersection profiles either in the form of a triangle or a quadrilateral, are computed, with the summation amount to the original total illumination area (which is normalized to 1).

We further accumulate the normalized area of each intersecting polygon over all mesh vertices that are located on the two ends of the intersecting edges, with a weight inversely proportional to the edge-segment lengths between the node and the edge-intersection. After this operation, the total illumination intensity of 1 is distributed along all nodes immediately located above or below the virtual source plane, where the edges between the nodes intersect with the virtual source plane.


\section{Nodal and element-based adjoint method for computing the Jacobian matrix}
Equation (\ref{eq:weakformpartial}) for each element written in
matrix form is
\begin{equation}
\label{eq:matrix} \mathbf{A}{\boldsymbol\Phi}=\mathbf{b}
\end{equation}
where the element of $\mathbf{A}$ can be written as
\begin{equation}
a_{i,j}=\sum_{i,j\in e}\left(\sum_{k=1}^4D_k\left\langle
\phi_k(r)\nabla\varphi_i(r)\cdot\nabla\varphi_j(r)\right\rangle+\sum_{k=1}^4\left(\mu_{a}^k+\frac{j\omega}{c}\right)\left\langle\phi_k(r)\varphi_i(r)\varphi_j(r)\right\rangle\right)
\end{equation}
for node-based FEM matrix, and 
\begin{equation}
a_{i,j}=\sum_{i,j\in e}\left(D_e\left\langle
\nabla\varphi_i(r)\cdot\nabla\varphi_j(r)\right\rangle+\left(\mu_{a}^{e}+\frac{j\omega}{c}\right)\left\langle\varphi_i(r)\varphi_j(r)\right\rangle\right)
\end{equation}
for the element-based FEM matrix, where $D_e$ and $\mu_a^{e}$ are the diffusion and absorption coefficients 
of the $e$-th element.

Taking the derivative of $\mu_a$ and $D$ on both sides of
(\ref{eq:matrix}), we get
\begin{eqnarray}
\frac{\partial \mathbf{A}}{\partial
\mu_a^k}{\boldsymbol\Phi}&=&-\mathbf{A}\frac{\partial {\boldsymbol\Phi}}{\partial \mu_a^k}\\
\frac{\partial \mathbf{A}}{\partial
D_k}{\boldsymbol\Phi}&=&-\mathbf{A}\frac{\boldsymbol\Phi}{\partial
D^k}
\end{eqnarray}
where the elements of matrix $\frac{\partial \mathbf{A}}{\partial
\mu_a^k}$ and $\frac{\partial \mathbf{A}}{\partial D_k}$ are
written as
\begin{eqnarray}
k_{i,j}=\frac{\partial
a_{i,j}}{\mu_a^k}&=&\left\langle\phi_k(r)\varphi_i(r)\varphi_j(r)\right\rangle\\
h_{i,j}=\frac{\partial a_{i,j}}{D_k}&=&\left\langle
\phi_k(r)\nabla\varphi_i(r)\cdot\nabla\varphi_j(r)\right\rangle
\end{eqnarray}
respectively.

Similarly, taking derivatives with respect to $D_e$ and $\mu_a^{e}$ yields
\begin{eqnarray}
k_{i,j}=\frac{\partial
a_{i,j}}{\mu_a^{e}}&=&\left\langle\varphi_i(r)\varphi_j(r)\right\rangle\\
h_{i,j}=\frac{\partial a_{i,j}}{D_e}&=&\left\langle
\nabla\varphi_i(r)\cdot\nabla\varphi_j(r)\right\rangle
\end{eqnarray}
respectively. Again, the angular bracket $\langle\cdot\rangle$ represents volume integration inside
an element, i.e.
\begin{equation}
\langle f(r) \rangle=\int_{\Omega_e}f(r)dr
\end{equation}
Therefore, the deriatives to $\mu_a^e$ can be calculated using this relationship:
\begin{equation}
\int_{\Omega_e}\varphi_1^l\varphi_2^m\varphi_3^n\varphi_4^sdr=\frac{l!m!n!s!}{(l+m+n+s+3)!}6V_e
\end{equation}
As a result, for element-based $\mu_a$ Jacobian, we have
\begin{equation}
\mathbf{K}_e=\left[\frac{\partial a_{i,j}}{\mu_a^{e}}\right]=\sum_e\frac{V_e}{20}\left(
\begin{array}{llll}
2 & 1 & 1 & 1 \\
1 & 2 & 1 & 1 \\
1 & 1 & 2 & 1 \\
1 & 1 & 1 & 2 
\end{array}
\right)
\end{equation}
For element-based $D$ Jacobian, we have
\begin{eqnarray}\nonumber
\mathbf{H}_e&=&\left[\frac{\partial
a_{i,j}}{D^{e}}\right]\\&=&\sum_e\frac{1}{(6V_e)^2}\left(
\begin{array}{llll}
\mathbf{C}_1\cdot\mathbf{C}_1 & \mathbf{C}_1\cdot\mathbf{C}_2 & \mathbf{C}_1\cdot\mathbf{C}_2 & \mathbf{C}_1\cdot\mathbf{C}_4\\
\mathbf{C}_2\cdot\mathbf{C}_1 & \mathbf{C}_2\cdot\mathbf{C}_2 & \mathbf{C}_2\cdot\mathbf{C}_2 & \mathbf{C}_2\cdot\mathbf{C}_4\\
\mathbf{C}_3\cdot\mathbf{C}_1 & \mathbf{C}_3\cdot\mathbf{C}_2 & \mathbf{C}_3\cdot\mathbf{C}_2 & \mathbf{C}_3\cdot\mathbf{C}_4\\
\mathbf{C}_4\cdot\mathbf{C}_1 & \mathbf{C}_4\cdot\mathbf{C}_2 & \mathbf{C}_4\cdot\mathbf{C}_2 & \mathbf{C}_5\cdot\mathbf{C}_4
\end{array}
\right)
\end{eqnarray}
where $\mathbf{C}_i = [a_i,b_i,c_i] (i=1,4)$ are the linear coefficients in the expressions of the basis functions $\varphi_i$ as
\begin{equation}
\left(
\begin{array}{l}
\varphi_1\\
\varphi_2\\
\varphi_3\\
\varphi_4
\end{array}\right)=\frac{1}{6V_e}\left(
\begin{array}{llll}
6V_{01} & a_1 & b_1 & c_1 \\
6V_{02} & a_2 & b_2 & c_2 \\
6V_{03} & a_3 & b_3 & c_3 \\
6V_{04} & a_4 & b_4 & c_4 
\end{array}
\right)\left(
\begin{array}{l}
1\\
x\\
y\\
z
\end{array}\right)
\end{equation}
Here, the matrix is the inversion of the Jacobian ($\mathbf{J}_e$) of the tetrahedron, i.e.
\begin{equation}\frac{1}{6V_e}
\left(
\begin{array}{llll}
6V_{01} & a_1 & b_1 & c_1 \\
6V_{02} & a_2 & b_2 & c_2 \\
6V_{03} & a_3 & b_3 & c_3 \\
6V_{04} & a_4 & b_4 & c_4 
\end{array}
\right)=\textrm{inv}\left(
\begin{array}{llll}
1 & 1 & 1 & 1 \\
x_1 & x_2 & x_3 & x_4 \\
y_1 & y_2 & y_3 & y_4 \\
z_1 & z_2 & z_3 & z_4  
\end{array}
\right)=\textrm{inv}(\mathbf{J}_e)
\end{equation}
and the element volume $V_e=\textrm{det}(\mathbf{J}_e)/6$

For the Jacobian corresponding to the $\mu_a^e$ in the $e$-th element, the 
element-based Jacobian can be expressed as
\begin{eqnarray}\nonumber
J_{\mu_a}((s,r),e)&=&\frac{\partial
\Phi_{s,r}}{d\mu_a^{e}}\\&=&\left(\mathbf{K}_e{\boldsymbol\Phi}_s^e\right)^T{\boldsymbol\Phi}_r^e
\end{eqnarray}
where vectors $\boldsymbol{\Phi}_s^e$ and $\boldsymbol{\Phi}_r^e$
are defined by $\boldsymbol{\Phi}_s^e=\{\Phi_s(\vec{p}_k)\}$,
$\boldsymbol{\Phi}_r^e=\{\Phi_r(\vec{p}_k)\}, (k=1,2,3,4)$,
respectively.

For the Jacobian corresponding to the $\tau$-th $\mu_a$, the nodal
adjoint method was applied and the Jacobian for the measurement at
the $r$-th detector with illumination of the $s$-th source is
expressed as
\begin{eqnarray}\nonumber
J_{\mu_a}((s,r),\tau)&=&\frac{\partial
\Phi_{s,r}}{d\mu_a^{\tau}}\\&=&\sum_{n\in\Omega_\tau}\left(\frac{\sum_{e\in\Omega_n}
V_e}{4}\right)\phi(\vec{p}_n)\Phi_s(\vec{p}_n)\Phi_r(\vec{p}_n)
\end{eqnarray}
where $\sum_{n\in\Omega_\tau}$ refers to the summation over the
forward nodes which fall inside $\Omega_\tau$ and
$\sum_{e\in\Omega_n}$ refers to the summation over the forward
elements that share the n-th forward node; $V_e$ is the volume of
the element, $\Phi_s(\vec{p}_n)$ and $\Phi_r(\vec{p}_n)$ are the
forward field and the adjoint field at the n-th forward node,
respectively.

The Jacobian with respect to $D$ is computed by the traditional
element-based adjoint method, which is written as

\begin{eqnarray}\nonumber
J_{D}((s,r),\tau)&=&\frac{\partial
\Phi_{s,r}}{dD_{\tau}}\\&=&\sum_{e\in\Omega_\tau}\left(\mathbf{H}_\tau^e{\boldsymbol\Phi}_s^e\right)^T{\boldsymbol\Phi}_r^e
\end{eqnarray}
where $\sum_{e\in\Omega_\tau}$ denotes the summation over all
elements that in the immediate vicinity of the parameter node $\tau$;
matrix $\mathbf{H}_\tau^e$ has form of
\begin{equation}
h^\tau_{i,j}=\left\langle
\phi_\tau(r)\nabla\varphi_i(r)\cdot\nabla\varphi_j(r)\right\rangle
\end{equation}

\section{Gauss-Newton iterative image reconstruction}
As we derived above, the Jacobian matrix is simply the first order Fr\'echet derivative of a multi-input (parameters) and multi-output (measurements) function. The 1st order derivative also gives the gradient direction, using which, one can solve for the inverse problem more efficiently.

Overall, the goal of the image reconstruction problem is to solve the below optimization problem
\begin{equation}\label{eq:gw}
\underset{\mu}{\arg\min} \, || \mathbf{y-\Phi(\mu)} ||^2_2
\end{equation}
where $\mathbf y$ denote the measurements collected from all source-detector pairs, $\mathbf\Phi$ is the forward model that maps the parameters $\mu$ to the forward solution. In nearly all cases, the measurement $\mathbf y$ is not a complete sample of the forward solution, but rather a sparse subset acquired at discrete locations. To match the two terms, the second term in the above expression is often represented by $\mathbf{\Lambda\Phi(\mu)}$ where $\mathbf\Lambda$ is a sampling matrix. For simplicity, here we assume that $\mathbf\Lambda$ is absorbed in the forward operator $\mathbf\Phi$. We note here that the model output $\mathbf\Phi(\mu)$ is a nonlinear function of the parameter $\mu$. Thus, in order to solve for $\mu$ that minimizes (\ref{eq:gw}), iterative estimation methods based on numerical gradient are often used.

The Gauss-Newton method is one of the widely applied algorithms to solve the above optimization problem. It is a well-balanced algorithm with a convergence speed between the linear convergence algorithm, such as conjugate-gradient method, and the 2nd-order convergence algorithm, such as the Raphson-Newton algorithm. Compared to the Raphson-Newton algorithm, where the 2nd-order derivative, i.e. the Hessian matrix ($H=\frac{\partial^2 \mathbf\Phi}{\partial\mu^2}$), has to be computed, the Gauss-Newton method uses an approximated Hessian matrix, i.e. the Gauss-Newton Hessian matrix, defined by $G=J^TJ$ where $J$ is the Jacobian, i.e. the first order derivative.


\section{Data calibration and simultaneous source-detector coupling coefficient estimation}
The source-detector coupling coefficients (SD) are multiplicative factors associated with each source and detector, representing the optode-dependent variations, including varied source power, detector sensitivities, fiber coupling efficiency, optical fiber light leakage, bending, and fiber-tissue coupling (such as hair, pressure differences, etc.). In practice, such coupling coefficients are typically considered to be removed after applying the below calibration formula:
\begin{equation}\label{eq:calib}
\mathbf{\Phi}^{t}_{calib}=\frac{\hat{\mathbf{\Phi}}^{t}_{meas}}{\hat{\mathbf{\Phi}}^{p}_{meas}}\mathbf{\Phi}_{calc}^{p}
\end{equation}
If we assume both the target measurement $\hat{\mathbf{\Phi}}^{t}$ and the phantom measurement $\hat{\mathbf{\Phi}}^{p}$ both contain the multiplicative SD coefficients, $s_id_j$, for the $(i,j)$-th measurement pair, i.e. $\hat{\mathbf{\Phi}}^{t}=SD\times\mathbf{\Phi}^{t}$ and $\hat{\mathbf{\Phi}}^{p}=SD'\times\mathbf{\Phi}^{p}$. If we make the assumption that $SD=SD'$, applying Eq. \ref{eq:calib}, $sd$ will be removed from the measurement data.

However, in realistic measurements, $SD'$ may not be the same as $SD$, as a result, the calibrated data $\mathbf{\Phi}^{t}_{calib}$ may still present uncalibrated coupling differences. The uncalibrated $SD$ coefficients are unknown. It is possible to simultaneously fit for the optical properties using the measurement data as well as to fit for these unknown $SD$ coefficients.

If we denote $\mathbf{\Phi}'=sd\times\mathbf{\Phi}$, where $sd=s\otimes d$ is the Kronecker product of the per-source-channel coupling coefficients $s_i$ and the per-detector-channel coupling coefficients $d_j$. The simultaneous $sd$ estimation aims to solve the following equation in least-square estimation
\begin{equation}
\left(\mathbf{J}_{\mu_a},\mathbf{J}_{D},\mathbf{J}_{sd}\right)\left(
\begin{array}{c}
\boldsymbol{\Delta}{\mu_a}\\
\boldsymbol{\Delta}{D} \\
s\\d
\end{array}
\right)=\boldsymbol{\Phi}^{meas}-\boldsymbol{\Phi}^{calc}
\end{equation}
where the $sd$ Jacobian is defined as
\begin{equation}\label{eq:Jsd}
\mathbf{J}_{sd}=\left[\mathbf{J}_{s}, \mathbf{J}_{d}\right]
\end{equation}
and
\begin{eqnarray}\label{eq:JsJd}
\mathbf{J}_{s} &=& \frac{\partial(\{s_id_j\times\Phi(i,j)\})}{\{s_i\}}=\{d_j\Phi(i,j)\}\\
\mathbf{J}_{d} &=& \frac{\partial(\{s_id_j\times\Phi(i,j)\})}{\{d_j\}}=\{s_i\Phi(i,j)\}\\
\end{eqnarray}
In frequency-domain (FD) systems, the measurement $\mathbf\Phi$ are complex numbers, therefore, $\mathbf{J}_{sd}$ is also complex. In this case, the $s$ and $d$ coupling coefficients can also be complex numbers. This makes them different from $\Delta\mu_a$ and $\Delta_D$ which can only take real-values.

\section{Log-magnitude/unwrapped phase form of Jacobian}
The complex form update equation is
\begin{equation}
\left(\mathbf{J}_{\mu_a},\mathbf{J}_{D}\right)\left(
\begin{array}{l}
\boldsymbol{\Delta}{\mu_a}\\
\boldsymbol{\Delta}{D}
\end{array}
\right)=\boldsymbol{\Phi}^{meas}-\boldsymbol{\Phi}^{calc}
\end{equation}
or
\begin{equation}
\label{eq:cpxupdate} \mathbf{J}\left(
\begin{array}{l}
\boldsymbol{\Delta}{\mu_a}\\
\boldsymbol{\Delta}{D}
\end{array}
\right)=\boldsymbol{\Delta}\boldsymbol{\Phi}
\end{equation}
For FD-systems, the misfit vector $\mathbf{\Delta\Phi}$ is a complex number, so are the Jacobian matrices. In these cases the above equation can be solved in the real-form of (\ref{eq:cpxupdate}) is given by

\begin{equation}\label{eq:singlewavelength}
\left(
\begin{array}{ll}
\Re e{\mathbf{J}}_{\mu_a} & \Re e{\mathbf{J}}_{D}\\
\Im m{\mathbf{J}}_{\mu_a} & \Im m{\mathbf{J}}_{D}
\end{array}\right)
\left(
\begin{array}{l}
\boldsymbol{\Delta}{\mu_a}\\
\boldsymbol{\Delta}{D}
\end{array}
\right)=\left(
\begin{array}{l}
\Re e{\boldsymbol{\Delta\Phi}}\\
\Im m{\boldsymbol{\Delta\Phi}}
\end{array}\right)
\end{equation}
When $sd$ is used, the above equation can be further expanded as
\begin{equation}
\label{eq:realupdate}\left(
\begin{array}{lllr}
\Re e{\mathbf{J}}_{\mu_a} & \Re e{\mathbf{J}}_{D} & \Re e{\mathbf{J}}_{sd} & -\Im m{\mathbf{J}}_{sd}\\
\Im m{\mathbf{J}}_{\mu_a} & \Im m{\mathbf{J}}_{D} & \Im m{\mathbf{J}}_{sd} & \Re e{\mathbf{J}}_{sd}
\end{array}\right)
\left(
\begin{array}{l}
\boldsymbol{\Delta}{\mu_a}\\
\boldsymbol{\Delta}{D}\\
\Re e({sd})\\
\Im m({sd})
\end{array}
\right)=\left(
\begin{array}{l}
\Re e{\boldsymbol{\Delta\Phi}}\\
\Im m{\boldsymbol{\Delta\Phi}}
\end{array}\right)
\end{equation}

When the log-amplitude/unwrapped phase reconstruction~\cite{FangThesis, Fang2006a, Meaney2007b} is enabled, instead of fitting for 
the complex-valued misfit vector $\boldsymbol{\Delta\Phi}$, one fits the concatenated log-amplitude ($\boldsymbol{\Delta\Gamma(\Phi)}=\log(|\boldsymbol{\Phi}^{meas}|)-\log(|\boldsymbol{\Phi}^{calc}|)$) and the phase ($\boldsymbol{\Delta\angle(\Phi)}$)
vectors of the misfit. The above equation can be rewrite to

\begin{equation}\label{eq:logphase}
\left(
\begin{array}{lll}
\mathbf{J}_1 & \mathbf{J}_3 & \mathbf{J}_5\\
\mathbf{J}_2 & \mathbf{J}_4 & \mathbf{J}_6
\end{array}\right)
\left(
\begin{array}{l}
\boldsymbol{\Delta}{\mu_a}\\
\boldsymbol{\Delta}{D}\\
\Re e({sd})\\
\Im m({sd})
\end{array}
\right)=\left(
\begin{array}{l}
{\boldsymbol{\Delta\Gamma(\Phi)}}\\
{\boldsymbol{\Delta\angle(\Phi)}}
\end{array}\right)
\end{equation}

where
\begin{eqnarray}\nonumber
\mathbf{J}_1&=&\frac{\Re e{\boldsymbol{\Delta\Phi}}\Re e{\boldsymbol{\mathbf{J}_{\mu_a}}}
+\Im m{\boldsymbol{\Delta\Phi}}\Im m{\boldsymbol{\mathbf{J}_{\mu_a}}} }{|\boldsymbol{\Delta\Phi}|^2} \\\nonumber
\mathbf{J}_2&=&\frac{\Re e{\boldsymbol{\Delta\Phi}}\Im m{\boldsymbol{\mathbf{J}_{\mu_a}}}
-\Im m{\boldsymbol{\Delta\Phi}}\Re e{\boldsymbol{\mathbf{J}_{\mu_a}}} }{|\boldsymbol{\Delta\Phi}|^2} \\
\mathbf{J}_3&=&\frac{\Re e{\boldsymbol{\Delta\Phi}}\Re e{\boldsymbol{\mathbf{J}_{D}}}
+\Im m{\boldsymbol{\Delta\Phi}}\Im m{\boldsymbol{\mathbf{J}_{D}}} }{|\boldsymbol{\Delta\Phi}|^2} \\\nonumber
\mathbf{J}_4&=&\frac{\Re e{\boldsymbol{\Delta\Phi}}\Im m{\boldsymbol{\mathbf{J}_{D}}}
-\Im m{\boldsymbol{\Delta\Phi}}\Re e{\boldsymbol{\mathbf{J}_{D}}} }{|\boldsymbol{\Delta\Phi}|^2} \\\nonumber
\mathbf{J}_5&=& \frac{1}{|sd|} \\\nonumber
\mathbf{J}_6&=& I \nonumber
\end{eqnarray}
where $J_5$ and $J_6$ are approximate versions of the $SD$ Jacobian in this case.

\section{Multi-spectral image reconstruction}
The multi-spectral image reconstruction algorithm was initially developed by Dr. Fang in the context of microwave tomography~\cite{Fang2004a, FangThesis}. Dr. Fang's multi-spectral algorithm had inspired Shuba Srinivasan, a PhD student at the time working under the supervision of Dr. Brian Pogue in the same Department. A year later, Srinivasan \emph{et al.} ported this algorithm to NIR imaging and achieved simultaneous chromophore reconstruction using multiple wavelengths~\cite{Srinivasan2025} (unfortunately, the failed cite the original multi-spectral algorithm paper~\cite{Fang2004a}).

When data from multiple wavelengths are used simultaneously, one can reconstruct wavelength-independent chromorphore molar-concentrations $c_i$ or volume fractions $f_i$ instead of $\mu_a$ and $D$. The relationship between $\mu_a$ and $c$ is expressed as
\begin{equation}
\mu_a(\lambda)=\epsilon_{HbO}(\lambda)c_{HbO}+\epsilon_{HbR}(\lambda)c_{HbR}+\epsilon_{water}(\lambda)f_{water}+\epsilon_{lipids}(\lambda)f_{lipids}+ \cdots
\end{equation}
where $\lambda$ is the wavelength, $\epsilon_i(\lambda)$ are the extinction coefficients of each chromorphore, and $f_{water}$ and $f_{lipids}$ are the volume fractions of water and lipids in the tissue. The extinction coefficients over the wavelength can be found in the literature. 

The inverse-power law is used to express the reduced scattering coefficient $\mu_s'$ in terms of wavelengths, i.e.
\begin{equation}
\mu_s'(\lambda)=a\times\lambda^{-b}
\end{equation}
where $a$ is the scattering amplitude and $b$ is the scattering power. In some literature, a normalization is performed, making it easier for unit conversions
\begin{equation}
\mu_s'(\lambda)=a\times\left(\frac{\lambda}{500 \mathrm{nm}}\right)^{-b}.
\end{equation}
in this case, $a$ has the unit of $\mu_s'$ (1/mm) and $b$ becomes unit-less.

In the multi-spectral reconstructions, the following equation is solved
\begin{equation}\label{eq:multispectral}
\left(
\begin{array}{llll}
\mathbf{J}_{c}(\lambda_1) & \mathbf{J}_{a}(\lambda_1) & \mathbf{J}_b(\lambda_1) & \mathbf{J}_{sd}(\lambda_1)\\
\mathbf{J}_{c}(\lambda_2) & \mathbf{J}_{a}(\lambda_2) & \mathbf{J}_b(\lambda_2) & \mathbf{J}_{sd}(\lambda_2)\\
\mathbf{J}_{c}(\lambda_3) & \mathbf{J}_{a}(\lambda_3) & \mathbf{J}_b(\lambda_3) & \mathbf{J}_{sd}(\lambda_3)\\
 & \cdots & \cdots 
\end{array}\right)
\left(
\begin{array}{c}
\boldsymbol{\Delta}{c_{HbO}}\\
\boldsymbol{\Delta}{c_{HbR}}\\
\boldsymbol{\Delta}{f_{water}}\\
\boldsymbol{\Delta}{f_{lipids}}\\
a\\
b\\
\Re e({sd})\\
\Im m({sd})
\end{array}
\right)=\left(
\begin{array}{c}
\boldsymbol{\Delta\Phi(\lambda_1)}\\
\boldsymbol{\Delta\Phi(\lambda_2)}\\
\boldsymbol{\Delta\Phi(\lambda_3)}\\
\cdots
\end{array}\right)
\end{equation}
and $\mathbf{J}_{c}=\left[\mathbf{J}_{HbO}, \mathbf{J}_{HbR}, \mathbf{J}_{water}, \mathbf{J}_{lipids}, \cdots\right]$, with
\begin{equation}
\mathbf{J}_{c}^i(\lambda_j)=\epsilon_i(\lambda_j)\mathbf{J}_{\mu_a}(\lambda_j)
\end{equation}
$i=$\{HbO, HbR, water, lipids, ...\}; also, we have
\begin{eqnarray}
\mathbf{J}_a(\lambda)&=&-\frac{\lambda^{b}}{3a^2} = -\left[3D^2(\lambda)\lambda^{-b}\right]\mathbf{J}_D(\lambda)\\
\mathbf{J}_b(\lambda)&=& \frac{3a\lambda^{-b}}{\log(\lambda)} =\left[D(\lambda)\log(\lambda)\right]\mathbf{J_D(\lambda)}
\end{eqnarray}
Here we assume that the $sd$ coefficients are wavelength-independent.

\section{Solving the inverse problem}
The perturbation equation in the 1st order, i.e. Eqs. \ref{eq:multispectral}, \ref{eq:singlewavelength} and \ref{eq:logphase} can be solved using ordinary least-square (OLS), weighted least-square (WLS) or generalized least-square (GLS) methods~\cite{FangThesis}.

The OLS solution solves the below parameter estimation problem
\begin{equation}
\underset{\mu}{\arg\min} \, || \mathbf{y-\Phi(\mu)} ||^2_2
\end{equation}
Assuming iid Gaussian noise in the measurement $\mathbf y$, the optimal solution is given by
\begin{eqnarray}
\mathbf{\Delta\mu}&=&(\mathbf{J^TJ}+\lambda I)^{-1}\mathbf{J}^T\mathbf{\Delta\Phi},\;\mathrm{or}\label{eq:overols}\\
\mathbf{\Delta\mu}&=&\mathbf{J}^T(\mathbf{JJ}^T+\lambda I)^{-1}\mathbf{\Delta\Phi}\label{eq:underols}
\end{eqnarray}
where $\lambda$ denotes the Tikhonov regularization parameter. Both of the above forms are equivalent, according to \emph{Sherm-Morrison-Woodbury
identity}. However, due to the matrix sizes, it is generally preferred to solve (\ref{eq:overols}) when $\mathbf{J}$ has more rows than columns (i.e. overdetermined form: number of measurements is more than that of unknowns) and solve (\ref{eq:underols}) when $\mathbf{J}$ has more columns than rows (i.e. under-determined form: number of unknowns is more than that of the measurements).

If WLS is used, the measurement data are considered random variables that may have independent variances. The residual is weighted according to the inverse square-root of the covariance matrix ($\mathbf{C_y}$) of the data, so that the data presenting the highest variance will be trusted the least, and those presenting the lowest variance will contribute the most in the residual calculations. This is equivalent to solving the below optimization problem
\begin{equation}
\underset{\mu}{\arg\min} \, || \mathbf{C_y}^{-1/2}\left(\boldsymbol{y-\Phi(\mu)}\right) ||^2_2
\end{equation}
and this gives us the solution as
\begin{eqnarray}
\mathbf{\Delta\mu}&=&(\mathbf{J^TC_y^{-1}J}+\lambda I)^{-1}\mathbf{J^TC_y}^{-1}\mathbf{\Delta\Phi},\;\mathrm{or}\label{eq:overwls}\\
\mathbf{\Delta\mu}&=&\mathbf{J^T}(\mathbf{JJ^T}+\lambda \mathbf{C_y})^{-1}\mathbf{\Delta\Phi}\label{eq:underwls}
\end{eqnarray}

In the case of GLS, both the measurements and the unknowns are considered random variables, and their ``assumed'' covariance matrices, i.e. $\mathbf{C_y}$ and $\mathbf{C_\mu}$ respectively, as \emph{a priori} information, should be used. This is equivalent to the below optimization problem
\begin{equation}
\underset{\mu}{\arg\min} \, || \mathbf{C_y}^{-1/2}\left(\boldsymbol{y-\Phi(\mu)}\right) ||^2_2 + \lambda|| \mathbf{C_\mu}^{-1/2}\boldsymbol{\mu} ||^2_2
\end{equation}
and the corresponding update equation can be seen as
\begin{eqnarray}
\mathbf{\Delta\mu}&=&(\mathbf{J^TC_y^{-1}J}+\lambda \mathbf{C}_\mu^{-1})^{-1}\mathbf{J^TC_y}^{-1}\mathbf{\Delta\Phi},\;\mathrm{or}\label{eq:overgls}\\
\mathbf{\Delta\mu}&=&\mathbf{C_\mu J^T}(\mathbf{JC_\mu J^T}+\lambda \mathbf{C_y})^{-1}\mathbf{\Delta\Phi}\label{eq:undergls}
\end{eqnarray}
In reality, the covariance matrix of the measurements $\mathbf{C_y}$ can be estimated using repeated measurements, whereas the covariance matrix of the unknowns $\mathbf{C_\mu}$ need to be provided as prior information.

When the prior values (or the expectations) of the unknowns are assumed as $\boldsymbol\mu_0$, one can solve the below extended optimization problem as
\begin{equation}\label{eq:objfunglsprior}
\underset{\mu}{\arg\min} \, || \mathbf{C_y}^{-1/2}\left(\boldsymbol{y-\Phi(\mu)}\right) ||^2_2 + \lambda|| \mathbf{C_\mu}^{-1/2}(\boldsymbol{\mu-\mu_0}) ||^2_2
\end{equation}
and the corresponding solutions are
\begin{eqnarray}
\mathbf{\Delta\mu}&=&(\mathbf{J^TC_y^{-1}J}+\lambda \mathbf{C}_\mu^{-1})^{-1}\left[\mathbf{J^TC_y}^{-1}\mathbf{\Delta\Phi} - \boldsymbol{C_\mu^{-1}(\mu-\mu_0)}\right],\;\mathrm{or}\label{eq:overglsprior}\\\nonumber
\mathbf{\Delta\mu}&=&\left[ I - \mathbf{C_\mu J^T}(\mathbf{JC_\mu J^T}+\lambda \mathbf{C_y})^{-1}\mathbf{J}\right] \\
 && \times\left[\mathbf{C_\mu J^TC_y^{-1}\Delta\Phi} - \boldsymbol{(\mu - \mu_0)}\right]\label{eq:underglsprior}
\end{eqnarray}

\section{Tissue bulk optical property fitting}\label{sec:bulkfitting}
The tissue bulk optical property refers to the set of averaged optical property across a heterogeneous domain. The estimation of the tissue bulk properties is achieved via solving the same optimization problem as shown above, except that the unknown space is compressed dramatically and contains only a single set of parameters, one parameter per unknown specie. 

Using the general reconstruction equation Eq. \ref{eq:multispectral}, the bulk optical properties can be estimated by solving the below equation
\begin{equation}\label{eq:multispectralbulk}
\left(
\begin{array}{llll}
\sum_i{J}_{c}(\lambda_1) & \sum_i{J}_{a}(\lambda_1) & \sum_i{J}_b(\lambda_1) & \sum_i{J}_{sd}(\lambda_1)\\
\sum_i{J}_{c}(\lambda_2) & \sum_i{J}_{a}(\lambda_2) & \sum_i{J}_b(\lambda_2) & \sum_i{J}_{sd}(\lambda_2)\\
\sum_i{J}_{c}(\lambda_3) & \sum_i{J}_{a}(\lambda_3) & \sum_i{J}_b(\lambda_3) & \sum_i{J}_{sd}(\lambda_3)\\
 & \cdots & \cdots 
\end{array}\right)
\left(
\begin{array}{c}
\overline{\boldsymbol{\Delta}{c_{HbO}}}\\
\overline{\boldsymbol{\Delta}{c_{HbR}}}\\
\overline{\boldsymbol{\Delta}{f_{water}}}\\
\overline{\boldsymbol{\Delta}{f_{lipids}}}\\
\bar{a}\\
\bar{b}\\
\overline{\Re e({sd})}\\
\overline{\Im m({sd}})
\end{array}
\right)=\left(
\begin{array}{c}
\boldsymbol{\Delta\Phi(\lambda_1)}\\
\boldsymbol{\Delta\Phi(\lambda_2)}\\
\boldsymbol{\Delta\Phi(\lambda_3)}\\
\cdots
\end{array}\right)
\end{equation}
where $\sum_i$ denotes the summation of all the columns (the dimension corresponds to the unknowns) of the Jacobian matrices and $\overline{\boldsymbol{\Delta}{c_i}}$ is the update to correct the $i$-th bulk optical property from the initial background value.

The above optical property fitting algorithm can be generalized in the below two scenarios. In the first case, where a piecewise constant segmentation is provided (as priors), the above equation can be modified to fit multiple sets of optical property, one for each segmentation
\begin{equation}\label{eq:multispectralseg}
\left(
\begin{array}{llll}
\sum_{i\in\Omega_1}{J}(\lambda_1) & \sum_{i\in\Omega_2}{J}(\lambda_1) & \cdots & \sum_{i\in\Omega_{N_s}}{J}(\lambda_1) \\
\sum_{i\in\Omega_1}{J}(\lambda_2) & \sum_{i\in\Omega_2}{J}(\lambda_2) & \cdots & \sum_{i\in\Omega_{N_s}}{J}(\lambda_2) \\
\sum_{i\in\Omega_1}{J}(\lambda_3) & \sum_{i\in\Omega_2}{J}(\lambda_3) & \cdots & \sum_{i\in\Omega_{N_s}}{J}(\lambda_3) \\
 & \cdots & \cdots 
\end{array}\right)
\left(
\begin{array}{c}
\overline{\boldsymbol{\Delta}{\mu_{\Omega_1}}}\\
\overline{\boldsymbol{\Delta}{\mu_{\Omega_2}}}\\
\overline{\boldsymbol{\Delta}{\mu_{\Omega_3}}}\\
\cdots
\end{array}
\right)=\left(
\begin{array}{c}
\boldsymbol{\Delta\Phi(\lambda_1)}\\
\boldsymbol{\Delta\Phi(\lambda_2)}\\
\boldsymbol{\Delta\Phi(\lambda_3)}\\
\cdots
\end{array}\right)
\end{equation}
where $\Delta\mu_{\Omega_i}=\left[\overline{c_{HbO,\Omega_i}},\overline{c_{HbR,\Omega_i}},\cdots\right]^T$ is the average optical properties for the $i$-th segment. The above segmentation-based bulk-property fitting method is sometimes referred to as the ``hard-priors''.

If a probabilistic segmentation is provided to replace the binary segmentations, where every node is represented as a set of probabilities (or fraction of volumes) of being different tissue types, so that the probility map is defined as a matrix
\begin{equation}
P=\left\{p^k_{i}\right\}_{i=1,\cdots,N; k=1,\cdots,N_s}
\end{equation}
where $0\le p^k_{i} \le 1$ is the probability for the $i$-th node/voxel to be tissue type $k$, and $\sum_k{p^k_{i}}=1$ holds for any $i$. The above segmentation-based bulk-property fitting equation can be generalized to
\begin{equation}\label{eq:multispectralfuzzysegbulk}
\left(
\begin{array}{c}
\mathbf{J}(\lambda_1)P\\
\mathbf{J}(\lambda_2)P\\
\mathbf{J}(\lambda_3)P\\
\cdots 
\end{array}\right)
\left(
\begin{array}{c}
\overline{\boldsymbol{\Delta}{\mu_{\Omega_1}}}\\
\overline{\boldsymbol{\Delta}{\mu_{\Omega_2}}}\\
\overline{\boldsymbol{\Delta}{\mu_{\Omega_3}}}\\
\cdots
\end{array}
\right)=\left(
\begin{array}{c}
\boldsymbol{\Delta\Phi(\lambda_1)}\\
\boldsymbol{\Delta\Phi(\lambda_2)}\\
\boldsymbol{\Delta\Phi(\lambda_3)}\\
\cdots
\end{array}\right)
\end{equation}
Once the tissue-type averaged optical property update is calculated using the above equation, the node values are subsequently updated by $\{\Delta\mu_i\}_i=P\times\{\overline{\Delta\mu_{\Omega_k}}\}_{k=1,\cdots,N_s}$.

\section{Defining priors using the $L$-matrix}\label{sec:Lmatrix}
In many cases, Eq. \ref{eq:objfunglsprior} is solved in the following form in the context of Tikhonov regularization
\begin{equation}
\underset{\mu}{\arg\min} \, || \boldsymbol{y-\Phi(\mu)} ||^2_2 + \lambda|| L(\boldsymbol{\mu-\mu_0}) ||^2_2
\end{equation}
Compared with Eq. \ref{eq:objfunglsprior}, it is clear that the regularization matrix ($L$) is the assumed inverse square-root of the covariance matrix of unknowns ($\mathbf{C_\mu}^{-1/2}$), determined by the user as the \emph{a prior} information.

Various non-identity $L$-matrices have been proposed. For example, it can be defined as a simple Laplacian to penalize rapid changes between neighboring nodes or voxels. For an FEM-mesh based solution, $L$ can be defined as 
\begin{equation}
l_{i,j}=\left\{\begin{array}{cl}
-1 & \textrm{if $j$ and $j$ are connected and $i\ne j$}\\
0 &\textrm{if $i$ and $j$ are not connected and $i\ne j$}\\
N_i & \textrm{if $i=j$}\\
\end{array}\right.
\end{equation}
where $N_i$ is the degree of the node, denoting the total neighbor count of the node $i$.

The above Laplacian operator can be also applied to segmented domains derived from prior information (such as a co-registered structural-scan of the same volume), where the nodes/elements within the same tissue regions are labeled identically, we have
\begin{equation}
l_{i,j}=\left\{\begin{array}{cl}
-\frac{1}{N_k} & \textrm{if $i$ and $j$ belong to the same region $\Omega_k$}\\
0 &\textrm{if $i$ and $j$ belong to different regions}\\
1 & \textrm{if $i=j$}\\
\end{array}\right.
\end{equation}
where $N_k$ is the total node/voxels within the $k$-th region $\Omega_k$. The above method is sometimes referred to as the ``soft-priors''.

Another reported $L$-matrix example is the so-called ``Helmholtz operator''. 
\begin{equation}
l_{i,j}=\left\{\begin{array}{cl}
-\frac{1}{N_k +(h/h_f)^2}& \textrm{if $i$ and $j$ belong to the same region $\Omega_k$}\\
0 &\textrm{if $i$ and $j$ belong to different regions}\\
1 & \textrm{if $i=j$}\\
\end{array}\right.
\end{equation}
where $h$ is the average distance between nodes, and $h_f$ is the desired average image feature size.

\section{Compositional-prior guided reconstructions}

The ``compositional prior'' guided reconstruction algorithm was proposed in 2010~\cite{Fang2010}. In this algorithm, the $L$ matrix is also derived from a co-registered structure-image but with the added ability to fully utilize the detailed gray-scale image features~\cite{Deng2015a}. Different from all above mentioned prior-guided reconstruction algorithms, where spatial nodes are initially segmented as piecewise constant regions, the compositional prior utilizes probabilistic segmentation of the tissue, as described in Section \ref{sec:bulkfitting}. 

In short, the compositional prior approach penalizes reconstruct recovered property differences in the compositional space. It
computes the similarities in tissue compositions between spatial locations, derived from the intensity maps from \emph{a priori} structural images, and build a weighted graph connecting all nodes in the compositional space ($N_s$ dimensional space with $p^k$ as coordinates), then create the $L$-matrix as the Laplacian of the weighted graph in the compositional space.

Specifically, the $L$-matrix encodes the compositional prior is constructed as
\begin{equation}
l_{i,j}=\left\{\begin{array}{cl}
-\frac{u_{i,j}}{\beta\sqrt{d_id_j}}& \textrm{if $||C_i-C_j||_2<\alpha N_s$, i.e. connected}\\
0 &\textrm{if $||C_i-C_j||_2\ge\alpha N_s$, i.e. not connected}\\
1 & \textrm{if $i=j$}\\
\end{array}\right.
\end{equation}
where $||\cdot||_2$ is the $L_2$-norm, $N_s$ is the total number of tissue types, $C_i=\{p^1_i,p^2_i,\cdots,p^{N_s}_i\}$ is the compositional vector at $i$-th node, and $\alpha\in[0, 1]$ is a user-defined parameter to control the sparsity of the $L$-matrix - the lower the $\alpha$ value, the sparser the $L$-matrix; $\beta>1$ is another user-specified parameter to control the diagonal dominance of the $L$-matrix - the larger the $\beta$, the less penalization between nodes. The typical values for $\alpha$ is between 0.1 and 0.2, and that for $\beta$ is 1.2. 

Parameter $u_{i,j}$ is the weight between nodes $i$ and $j$, defined as the average per-component Euclidean distance in the compositional space between the two nodes, as $u_{i,j}=\alpha-||C_i-C_j||_2/N_s\ge 0$, and $d_i$ and $d_j$ are the ``degrees'' of the two nodes in a weighted graph, defined as $d_i=\sum_{k,||C_i-C_k||_2<\alpha N_s}{u_{i,k}}$.

As one can see, the above definition of $L$-matrix does not require piecewise constant segmentations, and avoids the loses of the gray-scale features from the structural images. The probabilistic segmentation is also better suited for medical imaging because a small region in the biological tissue typically contains a mixture of various types of tissues. The tissue probability can be treated as the volume fractions of different tissue types inside the smallest spatial discretization unit, such as a voxel or a tetrahedron.

It is apparent that there are many ways to segment the \emph{a priori} structural image intensity maps to the probability maps in the compositional space. We call this step as the ``fuzzy segmentation''. The above proposed algorithm is quite general and can accommodate many fuzzy segmentation methods. A simple linear mapping was proposed in the original paper, and various alternative mapping methods, including those based on the Gaussian-mixture models and threshold-based algorithms were studied in Deng 2015. From the latter paper, we concluded that the the reconstructed images are not sensitive to the fuzzy segmentation algorithm; however, all of these methods showed significantly around $2\times$ errors in the reconstructed images.

As we show in Section \ref{sec:Lmatrix}, the $L$-matrix is equivalent to the $\mathbf{C_\mu}^{-1/2}$ term in GLS, and $\mathbf{C_\mu}=(L^TL)^{-1}$ . Replacing $\mathbf{C_\mu}$ to (\ref{eq:undergls}), we have
\begin{eqnarray}
\mathbf{\Delta\mu}&=& (\mathbf{J^TJ}+\lambda \mathbf{L^TL})^{-1}J^T\mathbf{\Delta\Phi},\,\textrm{ or}\\
\mathbf{\Delta\mu}&=&\mathbf{(L^TL)^{-1} J^T}\left[\mathbf{J(L^TL)^{-1}J^T}+\lambda I\right]^{-1}\mathbf{\Delta\Phi}
\end{eqnarray}

When multiple properties are defined on every node, the compositional-prior derived $L$-matrix is shared between multiple blocks of the parameters. Given the block structure as shown in (\ref{eq:multispectral}), we have the following equation for the over-determined form
\begin{equation}
\mathbf{\Delta\mu}=
\left[\left(\begin{array}{l}
\mathbf{J}_1\\
\mathbf{J}_2\\
\mathbf{J}_3\\
\cdots 
\end{array}\right)
[\mathbf{J_1, J_2},\cdots]+\lambda
\left(\begin{array}{l}
\mathbf{C_\mu}\mathbf{J}_1\\
\mathbf{C_\mu}\mathbf{J}_2\\
\mathbf{C_\mu}\mathbf{J}_3\\
\cdots 
\end{array}\right)\right]^{-1}\mathbf{\Delta\Phi}
\end{equation}
Similarly, for the under-determined form, we have
\begin{eqnarray}\label{eq:underLTL}
\mathbf{\Delta\mu}&=&
\left(\begin{array}{l}
\mathbf{C_\mu}\mathbf{J}_1\\
\mathbf{C_\mu}\mathbf{J}_2\\
\mathbf{C_\mu}\mathbf{J}_3\\
\cdots 
\end{array}\right)\\\nonumber
&\times& \left[
\left(\begin{array}{cccc}
\mathbf{J}_1\mathbf{C_\mu}\mathbf{J}^T_1 & 0 & 0 & \cdots\\
0& \mathbf{J}_2\mathbf{C_\mu}\mathbf{J}^T_2 & 0 & \cdots\\
0 & 0 & \mathbf{J}_3\mathbf{C_\mu}\mathbf{J}^T_3 & \cdots\\
& \cdots& 
\end{array}\right)
+\lambda I\right]^{-1}\mathbf{\Delta\Phi}
\end{eqnarray}
where the Jacobian blocks $\mathbf{J}_k$ are the Jacobians corresponding to each optical property (chromorphore concentrations, absorption, scattering etc.) defined on the nodes.

In Deng 2015, we propose a fast algorithm to compute the above under-determined form by first computing the QR decomposition of $L$ as $L=QR$, and then compute $R^{-1}$ and construct an intermediate matrix $Z_i=\mathbf{J}_iR^{-1}$. Therefore, we have $\mathbf{C_\mu}\mathbf{J}_i=R^{-1}Z^T$ and $\mathbf{J}_i\mathbf{C_\mu}\mathbf{J}^T_i=Z_iZ^T$. 
Expanding these two terms in (\ref{eq:underLTL}), we can solve for the under-determined solution more efficiently.

\section{Advanced reconstruction techniques}

To be expanded.

\subsection{Dual-mesh method for solving decoupled forward and inverse problems}

\subsection{Multiple right-hand-side linear solvers for accelerating matrix solutions}

\subsection{Joint use of continuous-wave and frequency-domain data}

\subsection{Matrix scaling and variance equalization}

\subsection{Emperical Tikhonov regularization method}

\subsection{Two-dimensional prior guided 3D reconstructions}

\subsection{Other supported linear and least-square solvers}

\chapter{Software implementations}

To be expanded.

\clearpage

\bibliographystyle{unsrt}
\addcontentsline{toc}{chapter}{Bibliography}
\bibliography{redbird_manual}

\end{document}
